% LEMMA B' RESULT.
% Corrected AS/GKZ Hodge polygon, genuine Frobenius factor, and binary gate.

\subsection{The eigenvalue Newton polygon of the marked coordinate}
\label{ssec:lemmaBprime-result}

We retain the marked Laurent polynomial of
Theorem~\ref{thm:oracleC-marked-coordinate} and the exact Jacobi skeleton of
Theorem~\ref{thm:lemmaA-skeleton}.  For $q=p^r$ and
$1\leq j\leq p-2$, put
\[
 A=\chi_{p,j}\circ N_{\mathbb F_q/\mathbb F_p},
 \qquad
 S_{p,j}^{(r)}=-\sum_{x\in U(\mathbb F_q)}A(\Lambda(x)).
\]
For an auxiliary prime $\ell\ne p$, write $E_\ell$ for a coefficient field
containing the values of $A$.
Thus $S_{p,j}^{(1)}=\widetilde c_{p,j}^{\rm tor}$ and
$S_{p,j}^{(1)}\equiv b_j\pmod{\mathfrak p}$.  We use the effective sign
convention in which $S_{p,j}^{(r)}$ is a Frobenius power trace.  The signed
trace on $R\Gamma_c(U,\mathcal L_{p,j})$ is its negative.  This sign will be
restored in the slope-zero inclusion below.

\subsubsection{The corrected Adolphson--Sperber object}

The coordinate change already used in the proof of
Proposition~\ref{prop:oracleC-uniform-complexity},
\[
 u=\frac{1+y}{y},\qquad v=\frac{1+z}{z},\qquad
 w=-\frac{x}{uv},
\]
gives
\begin{equation}
 \Lambda=\frac{uv(w-1)(1-uvw)}{w(u-1)(v-1)},
 \qquad
 U=(\mathbb G_m)^3\setminus\{u=1,v=1,w=1,uvw=1\}.
 \label{eq:lemmaBprime-arrangement}
\end{equation}

There is an important trap here.  Converting the Kummer sum to an additive
sum with the phase $t\Lambda$ does \emph{not} permit direct application of
the nondegenerate Newton-polytope theorem.  At $x=y=-1$ two factors of
$\Lambda$ vanish, so $t\Lambda$ and all four toric logarithmic derivatives
vanish simultaneously.  In fact, if $P=\operatorname{Newt}(\Lambda)$, then
\[
 P=\{(X,Y,Z):-1\leq X,Y,Z\leq1,\ X-Y\leq1,\ X-Z\leq1\}.
\]
Its cross-sectional area is $4$ for $-1\leq X\leq0$ and $(2-X)^2$ for
$0\leq X\leq1$, whence
\begin{equation}
 3!\operatorname{vol}(P)=6\left(4+\int_0^1(2-X)^2\,dX\right)=38.
 \label{eq:lemmaBprime-wrong-volume}
\end{equation}
The same number is the normalized volume of
$\operatorname{conv}(0,P\times\{1\})$.  A naive use of this polytope would
therefore manufacture $38$ slots.  It is incompatible with the cancelled
rank below (and with $\chi_c(U)=1-7+17-14=-3$), precisely because the phase
is degenerate.

The correct Adolphson--Sperber/GKZ object follows by first performing the
Jacobi convolution.  Base change of
Theorem~\ref{thm:lemmaA-skeleton} gives the exact identity
\begin{equation}
 S_{p,j}^{(r)}=-\frac1{q-1}\sum_\rho
 J_q(\bar\rho,A)^2J_q(\bar\rho,A\rho)^2,
 \label{eq:lemmaBprime-A1-extension}
\end{equation}
where \emph{all} $q-1$ characters of $\mathbb F_q^\times$ occur.  With
$g_q(C)=\sum_tC(t)\psi(t)$ and $g_q(\varepsilon)=-1$, including the
inverse-character degeneracies, one has
\begin{equation}
 J_q(\bar\rho,A)^2J_q(\bar\rho,A\rho)^2
 =q^{-2}g_q(\bar\rho)^4g_q(A\rho)^2g_q(A^{-1}\rho)^2.
 \label{eq:lemmaBprime-gauss-eight}
\end{equation}
For example, when $\rho=A$ the Jacobi factor
$J_q(A^{-1},A)^2$ and the product $g_q(A^{-1})^2g_q(A)^2=q^2$
give the same formula; no exceptional summand has been dropped.
To fix all signs, inversions, and Tate normalizations, define
\[
 \mathcal H_{q,A}(t):=-\frac{1}{q^2(q-1)}
 \sum_\rho g_q(\bar\rho)^4g_q(A\rho)^2
                    g_q(A^{-1}\rho)^2\rho(t).
\]
This is our trace convention for the Tate-normalized balanced
hypergeometric object
\begin{equation}
 \operatorname {Hyp}(\varepsilon,\varepsilon,\varepsilon,\varepsilon;
 A,A,A^{-1},A^{-1}).
 \label{eq:lemmaBprime-hyp}
\end{equation}
Equations~\eqref{eq:lemmaBprime-A1-extension} and
\eqref{eq:lemmaBprime-gauss-eight} say exactly that
$\mathcal H_{q,A}(1)=S_{p,j}^{(r)}$.  The opposite Katz ordering exchanges
the two parameter lists and replaces $t$ by $t^{-1}$, hence gives the same
stalk at $t=1$.
Introduce the usual hypergeometric parameter $t$.  Expanding the eight
Gauss factors gives the associated GKZ circuit, which is nondegenerate for
generic $t\ne1$; this, rather than the degenerate phase $t\Lambda$, is the
object to which the twisted-height recipe of
Adolphson--Sperber~\cite{AS1993} applies.  The value $t=1$ is the unique
balanced-circuit singularity and must subsequently be treated by middle
extension and nearby cycles.

Here is the requested closed form.  Put
\[
 a=\left\langle\frac{j}{p-1}\right\rangle\in(0,1),\qquad
 \boldsymbol\alpha=(0,0,0,0),
\]
and sort the fractional-part data
\begin{equation}
 \boldsymbol\beta(a)
 =\operatorname {sort}\bigl(
   \langle a\rangle,\langle a\rangle,
   \langle-a\rangle,\langle-a\rangle\bigr)
 =\operatorname {sort}(a,a,1-a,1-a).
 \label{eq:lemmaBprime-beta}
\end{equation}
For $k=1,\ldots,4$, the twisted lattice-height statistic is
\begin{equation}
 \rho_a(k)=\#\{i:\alpha_i<\beta_k(a)\}-k=4-k.
 \label{eq:lemmaBprime-rho}
\end{equation}
Indeed, every entry of~\eqref{eq:lemmaBprime-beta} is strictly positive,
including at $a=1/2$.  Thus the generic circuit has Hodge polynomial
\begin{equation}
 H_a^{\rm gen}(T)=\sum_{k=1}^4T^{\rho_a(k)}
 =1+T+T^2+T^3.
 \label{eq:lemmaBprime-generic-HP}
\end{equation}

The two character multisets in~\eqref{eq:lemmaBprime-hyp} are disjoint, so
the generic hypergeometric object is irreducible of rank four.  The balanced
hypergeometric local-monodromy theorem (in the convolution normalization of
Katz~\cite{Katz2012}) says that the monodromy at $t=1$ is a nontrivial
pseudoreflection.  Its special eigenvalue is
\[
 \exp\!\left(2\pi i\left(\sum_k\beta_k-
                  \sum_i\alpha_i\right)\right)=\exp(4\pi i)=1.
\]
Thus $T_1=\exp N$ with $N\ne0$, $N^2=0$, and $\operatorname {rank}N=1$.
The standard rank-one-unipotent hypergeometric nearby-cycle calculation now
applies.  The parameter multiset is invariant under complex conjugation, so
the variation has a real form.  The polarized monodromy filtration $W(N)[3]$
has rank-one graded pieces in weights $2$ and $4$; the real form forces them
to have types $(1,1)$ and $(2,2)$, respectively.  Together with the four
generic Hodge degrees in~\eqref{eq:lemmaBprime-generic-HP}, this gives, after
one complex coefficient embedding,
\[
 \operatorname {Gr}_2^W\text{ of type }(1,1),\qquad
 \operatorname {Gr}_3^W\text{ of types }(3,0)+(0,3),\qquad
 \operatorname {Gr}_4^W\text{ of type }(2,2),
\]
where $N:\operatorname {Gr}_4^W\buildrel\sim\over\longrightarrow
\operatorname {Gr}_2^W(-1)$.  The middle extension has stalk
$\ker(T_1-1)=\ker N=W_3$, and strictness computes its induced Hodge
filtration from the graded types above.
It therefore removes the slope-$2$ slot, not the slope-$0$ slot.  The toric Hodge polynomial
and the requested multiplicity are
\begin{equation}
 H_{a,t=1}(T)=1+T+T^3,\qquad
 \boxed{\ h^0(j)=[T^0]H_{a,t=1}(T)=1\ }
 \quad(1\leq j\leq p-2).
 \label{eq:lemmaBprime-h0}
\end{equation}
This is a closed fractional-part formula; it uses no Ap\'ery value and no
numerical input.

For completeness, the three slots can also be seen on the arrangement.
The only non-simple-normal-crossings point is $(u,v,w)=(1,1,1)$.  On its
exceptional plane, with local coordinates denoted by $a,b,c$, the leading
Kummer function is
\[
 \Lambda_E=-\frac{c(a+b+c)}{ab}.
\]
For every nontrivial character $A$ over $\mathbb F_q$, setting $b=1$ gives
\begin{align*}
 \sum_{E^\circ(\mathbb F_q)}A(\Lambda_E)
 &=A(-1)\sum_{a,c\ne0}A(c)A(a+1+c)A^{-1}(a)=q.
\end{align*}
For $c\ne-1$ the inner $a$-sum is $-1$, while at $c=-1$ it is $q-1$;
the last equality follows after summing $A(c)$.  In perverse normalization,
the Mellin identity identifies the $t=1$ stalk with the effective class
$-[R\Gamma_c]$.  The blow-up weight spectral sequence identifies the
exceptional-plane contribution with its unique rank-one weight-$2$ piece,
$\operatorname {Gr}_2^W(\ker N)\simeq E_\ell(-1)$ above, and the displayed
sum says that geometric Frobenius acts on it by $q=p^r$.  Hence its
eigenvalue in the effective polynomial is exactly $+p$; the outer minus sign
has already been absorbed in $-[R\Gamma_c]$.

The involution
\begin{equation}
 \tau(u,v,w)=\bigl(w,(uvw)^{-1},u\bigr),\qquad
 \Lambda\circ\tau=\Lambda^{-1},
 \label{eq:lemmaBprime-involution}
\end{equation}
is defined over $\mathbb F_p$ and carries $\mathcal L_A$ to
$\mathcal L_{A^{-1}}$.  Composing this identification with Poincar\'e
duality gives a Frobenius-equivariant alternating perfect pairing
\[
 V_A\otimes_{E_\ell}V_A\longrightarrow E_\ell(-3)
\]
on the remaining rank-two quotient, with no residual constant Kummer twist.
In rank two the determinant equals the similitude multiplier, hence it is
$p^3$.  Therefore the effective characteristic factor is
the theorem-level identity
\begin{equation}
 \boxed{
 P_{p,j}(X)=(X-p)\bigl(X^2-a_{p,j}X+p^3\bigr),
 \qquad
 a_{p,j}=S_{p,j}^{(1)}-p.}
 \label{eq:lemmaBprime-factor}
\end{equation}
The equality is an equality of the semisimplified effective virtual
Frobenius class $-[R\Gamma_c(U,\mathcal L_{p,j})]$ with the rank-three
middle-extension stalk.  Thus it does not confuse the $q-1$ Mellin
summands in~\eqref{eq:lemmaBprime-A1-extension} with eigenvalues.

\subsubsection{Newton polygons and the inclusion theorem}

Since $a_{p,j}\equiv S_{p,j}^{(1)}\equiv b_j\pmod{\mathfrak p}$, put
$m=v_{\mathfrak p}(a_{p,j})$.  The two roots of the quadratic factor in
\eqref{eq:lemmaBprime-factor} have valuations
\begin{equation}
 \begin{cases}
 (0,3),&m=0,\\
 (1,2),&m=1,\\
 (3/2,3/2),&m\geq2.
 \end{cases}
 \label{eq:lemmaBprime-quadratic-slopes}
\end{equation}
Adding the Tate slope $1$ gives
\begin{equation}
 m_0(p,j)=
 \begin{cases}
 1,&b_j\not\equiv0\pmod p,\\
 0,&b_j\equiv0\pmod p.
 \end{cases}
 \label{eq:lemmaBprime-unit-count}
\end{equation}
In particular, $m_0(p,j)$ is never at least two.  At an ordinary point the
Newton slopes are $(0,1,3)$ and equal the Hodge slopes in
\eqref{eq:lemmaBprime-h0}; at a Hasse zero the polygon rises strictly above
the Hodge polygon.

We record explicitly the theorem required in the specification.  In the
effective convention, the unique unit root reduces to $b_j$.  Restoring
the sign of $R\Gamma_c$, its signed slope-zero sum is therefore
$-b_j\pmod{\mathfrak p}$.  Consequently
\begin{theorem}[Slope-zero inclusion]
\label{thm:lemmaBprime-inclusion}
For every prime $p\geq5$,
\begin{equation}
 \mathcal Z_p\subseteq
 \{1\leq j\leq p-2:m_0(p,j)\ne1\}.
 \label{eq:lemmaBprime-inclusion}
\end{equation}
For this marked coordinate equality holds in
\eqref{eq:lemmaBprime-inclusion}: both sides are exactly the Hasse-zero set.
\end{theorem}

Thus the inclusion is a theorem, but it is also a tautology: the moving
unit root exists exactly when its Hasse invariant $b_j$ is nonzero.

\paragraph{Extension-field computation.}
To check~\eqref{eq:lemmaBprime-factor} without making the circular
Jacobi-unit count, the script evaluates the full base-changed skeleton
\eqref{eq:lemmaBprime-A1-extension}.  If $\gamma$ generates
$\mathbb F_q^\times$, set
\[
 f_n=A(1-\gamma^n),\qquad
 g_n=A\left(\frac{\gamma^n}{1+\gamma^n}\right).
\]
Multiplicative Fourier inversion turns the complete $q-1$ character sum
into the exact cyclic convolution
\begin{equation}
 S_{p,j}^{(r)}=-\sum_{c\in\mathbb Z/(q-1)}(f*f)_c(g*g)_c.
 \label{eq:lemmaBprime-cyclic-engine}
\end{equation}
This is the promised $\mathbb F_{p^r}$ use of
Theorem~\ref{thm:lemmaA-skeleton}.  Only the characters with index
$K=(q-1)k/(p-1)$ are Hasse--Davenport lifts of base-field characters; the
other extension characters cannot be omitted.  At $(p,j,r)=(5,1,2)$ the
correct trace is $-125$, whereas the invalid termwise lifting shortcut
gives $-153$.

Newton's identities recover the cubic coefficients from $S_1,S_2,S_3$:
\[
 e_1=S_1,\qquad e_2=\frac{S_1^2-S_2}{2},\qquad
 e_3=\frac{S_1^3-3S_1S_2+2S_3}{6}.
\]
On six rational probe pairs, the portable direct-count and full-$q-1$ A1
engines independently compute all three traces and give
$e_2=pS_1+p^3-p^2$ and $e_3=p^4$, as predicted by
\eqref{eq:lemmaBprime-factor}.  The $p=13$ row below uses an independently
computed $S_1$ and the theorem-level factor.  Representative rows are
\begin{center}
\begin{tabular}{c|c|r|r|r|c}
 $p$ & $j$ & $S_1$ & $e_2$ & $e_3$ & Newton slopes \\ \hline
 $5$  & $1$ & $-5$  & $75$  & $625$   & $(1,1,2)$ \\
 $5$  & $2$ & $3$   & $115$ & $625$   & $(0,1,3)$ \\
 $7$  & $1$ & $-2$  & $280$ & $2401$  & $(0,1,3)$ \\
 $7$  & $3$ & $31$  & $511$ & $2401$  & $(0,1,3)$ \\
 $11$ & $5$ & $-33$ & $847$ & $14641$ & $(1,1,2)$ \\
 $13$ & $1$ & $5$   & $2093$& $28561$ & $(0,1,3)$
\end{tabular}
\end{center}
For example, the complete traces at $(p,j)=(7,1)$ are
\[
 (S_1,S_2,S_3,S_4)=(-2,-556,8875,133128),
\]
and at $(11,5)$ they are $(-33,-605,91839)$.  A genuinely cyclotomic
check at $p=11,j=1$, with $\zeta=\zeta_{10}$, gives
\[
 S_1=8\zeta^3-8\zeta^2-27,\qquad
 e_2=88\zeta^3-88\zeta^2+913,\qquad e_3=14641.
\]

For every prime $p\leq60$ and every nontrivial $j$ ($405$ pairs), an
independent production scan computed the complete traces $S_1,S_2,S_3$
modulo $p^6$ from~\eqref{eq:lemmaBprime-cyclic-engine}.  Thus every one of
the $q-1$ extension characters was retained for $q=p,p^2,p^3$; neither the
Ap\'ery recurrence nor~\eqref{eq:lemmaBprime-factor} was used to construct
the traces.  The optional \texttt{--full-small-scan} mode of
\texttt{lemmaBprime\_explore.py}, run under \texttt{sage -python}, is the
reproduction path; the standard portable verifier uses the exact anchor
probes instead of repeating this seven-minute scan.  Newton's identities
then gave
\begin{center}
\begin{tabular}{c|r|r|r}
 range & $(0,1,3)=\mathrm{NP}=\mathrm{HP}$ & $(1,1,2)$ & $m_0\geq2$ \\ \hline
 $p\leq60$, all $j$ & $390$ & $15$ & $0$
\end{tabular}
\end{center}
The $15$ nonordinary points are
\[
\begin{gathered}
 (5;1,3),\ (11;5),\ (17;3,13),\ (19;8,10),\\
 (31;8,22),\ (37;17,19),\ (41;10,30),\ (59;9,49).
\end{gathered}
\]
They are exactly the $15$ pairs with $b_j=0$ in this range.

For every prime $p\leq300$, the second range run used the already proved factor
\eqref{eq:lemmaBprime-factor}: it evaluated the base-field A1 skeleton
$p$-adically at every $j\in\mathcal Z_p$, and classified $20$ deterministic
random coordinates per prime from the factor and their base trace modulo $p$
(all available coordinates for the smallest primes).
It contained $63$ zero coordinates and $1140$ sampled coordinates.  All $63$ zeros had
$v_{\mathfrak p}(a_{p,j})=1$ and slopes $(1,1,2)$.  The random sample split
as $1127$ ordinary coordinates with slopes $(0,1,3)$ and $13$ hits on the
already enumerated zero set, again with slopes $(1,1,2)$.  No pair had two
slope-zero roots.  Thus the $p\leq60$ table is a full extension-field check
of the factor, while the $p\leq300$ table is a theorem-driven enumeration
plus independent base-A1 valuation checks at all Hasse zeros.  Neither table
is used to prove the non-numerical Hodge formula
\eqref{eq:lemmaBprime-h0}.

\subsubsection{Exact eigenvalue tests and the discriminant}

The polynomial
\[
 \Delta(t)=t^2-34t+1
\]
does occur naturally, but not as a Frobenius eigenline.  From
\eqref{eq:lemmaBprime-arrangement}, the logarithmic critical equations are
\[
 u^2vw-2uvw+1=0,\qquad
 uv^2w-2uvw+1=0,\qquad
 uvw^2-1=0.
\]
On $U$ they give
\[
 u=v=2-w,\qquad w^2-2w-1=0,
\]
and substitution yields the two critical values
\begin{equation}
 \Lambda_{\rm crit}=17\pm12\sqrt2=(1\pm\sqrt2)^4.
 \label{eq:lemmaBprime-critical-values}
\end{equation}
Thus the roots of $\Delta$ are conifold critical values of the map
$\Lambda$.  The involution~\eqref{eq:lemmaBprime-involution} exchanges
them.  The actual rank-two quotient remains self-dual; neither critical
value thereby becomes a rank-one Teichm\"uller eigenline.

The exact tests make this distinction visible.
At the split prime $p=7$, the roots of $\Delta$ are $2$ and $4$.  For
$j=1,2,3$ the quadratic factors are
\[
 X^2+9X+343,\qquad X^2+25X+343,\qquad X^2-24X+343.
\]
Their unit-root residues are respectively $5,3,3$, whereas the proposed
Teichm\"uller targets are respectively $\{2,4\}$, $\{2,4\}$, and
$\{1\}$.  Hence the candidate already fails modulo $7$.  More strongly,
for $j=1,2$ reduction modulo $\Phi_3(X)=X^2+X+1$ leaves
$8X+342$ and $24X+342$, and for $j=3$ the last quadratic evaluates to
$320$ at $X=1$.  There is no exact cyclotomic match.

The three quadratic eigenvalue fields in this test are
\[
 \mathbb Q(\sqrt{-1291}),\qquad
 \mathbb Q(\sqrt{-83}),\qquad
 \mathbb Q(\sqrt{-199}),
\]
none of which is the base Jacobi field
$E_7=\mathbb Q(\zeta_6)=\mathbb Q(\sqrt{-3})$.  Thus no base-field
Jacobi-sum monomial, affine in $j$ or otherwise, equals the unit root in
these exact cases.  In $\mathbb Q_7$, the unit residues $5,3,3$ are all
nonsquares; the other two slopes are $1$ and $3$, both odd.  Hence none of
the three actual eigenvalues is a square in the natural coefficient
completion.  This rules out the nonvacuous local Sym$^2$ square test.

At the inert prime $p=11$, the roots of $\Delta$ lie in
$\mathbb F_{11^2}$ and have order $6$.  For $j=1,2,4$ their powers do not
even lie in $\mathbb F_{11}$; for $j=3$ the target is $-1$ but
$b_3\equiv4\pmod {11}$; and $j=5$ has no unit root.  Residue coincidences
can occur: at $j=6$ both the target and the unit-root residue are $1$.

Fix an embedding
$\iota_p:\overline{\mathbb Q}\hookrightarrow\overline{\mathbb Q}_p$.
A Teichm\"uller lift of either reduction of $17\pm12\sqrt2$ is
$\iota_p(\xi)$ for a prime-to-$p$ root of unity $\xi$; in the inert case it
lies in the unramified quadratic extension.  Every complex conjugate of
$\xi$ has absolute value $1$.  By contrast, each root of the pure
weight-$3$ quotient is a Weil $p$-number whose complex conjugates have
absolute value $p^{3/2}$.  Injectivity of $\iota_p$ therefore forbids exact
equality in $\overline{\mathbb Q}_p$ with
$\omega(17\pm12\sqrt2)^{\pm j}$, although isolated congruences modulo $p$
are possible.

The $p=7$ field and square computations are additional exact sample checks:
they exclude every base-field Jacobi monomial in those cases and every
square in the natural coefficient completion
$E_{7,\mathfrak p}=\mathbb Q_7$.  We do not claim a universal classification
of affine Jacobi monomials, or a square obstruction after arbitrary scalar
extension.  The sole universally identified rank-one eigenvalue is $p$, of
slope $1$, and is irrelevant to a slope-zero collision.

\paragraph{Binary verdict.}
Equation~\eqref{eq:lemmaBprime-h0} gives $h^0(j)=1$ for every nontrivial
twist, so Newton-above-Hodge allows at most one slope-zero eigenvalue.
Equation~\eqref{eq:lemmaBprime-unit-count} says that the Newton polygon
equals the Hodge polygon precisely away from the original Hasse zeros.
Thus the two-unit-root split required by G2 is structurally impossible,
independently of whether the single moving root has an isolated alternative
description.  The requested verdict is therefore
\begin{center}
 \fbox{\textbf{G1 (KILL).}}
\end{center}
The crystalline criterion is exactly the Hasse-invariant tautology
$p\mid b_j\Longleftrightarrow m_0(p,j)=0$.  There is no second unit root
with which a moving root could collide.  In particular,
the critical values $17\pm12\sqrt2$ do not open a unit-root
anti-concentration problem.

\paragraph{Fail-closed computational consistency verification.}
Running \texttt{python3 problems/3.2/lemmaBprime\_verify.py} ends with
\begingroup\scriptsize
\begin{verbatim}
PASS full A1 probes: 6 pairs r<=3; p13 r=1; cyclotomic p11,j1 S1
PASS cubic factor: 6 probed cubics; 5 r4 recurrences; p13 predicted
PASS p<=59 classification: 405 factor classifications; base-A1 v(a)=1 at 15/15 zeros
PASS p<=300 sample: base-A1 v(a)=1 at 63/63 zeros; sample 1140, zero hits 13
PASS candidate matches: critical values; p7 exact Teich/Jacobi/Q7-square rejection
PASS binary gate: G1 (h0=1; the only fixed branch is Tate slope 1)
SUMMARY: 11 PASS, 0 FAIL
\end{verbatim}
\endgroup
